\documentclass[11pt]{article}
\usepackage[margin=2.5cm]{geometry}
\usepackage{amsmath,booktabs,graphicx,hyperref}
\title{DSP avanzado y benchmark científico para perturbaciones de calidad de energía}
\author{Legacy}
\date{}
\begin{document}
\maketitle

\section{Alcance}
Este documento acompaña el paquete reproducible M1--M11. Los resultados numéricos
finales se incorporarán únicamente después de ejecutar el pipeline completo sobre
la versión congelada del dataset.

\section{Métodos}
Se comparan FFT, STFT, DWT, CWT y S-Transform con entradas idénticas. La selección
de parámetros se realiza sobre train/validation y el conjunto test se reserva para
evaluación.

\section{DWT}
Se evalúan db4, db6, db8, sym4, sym6, coif3 y coif5 con niveles válidos entre 3 y 7.
Se extraen energía por banda, entropía wavelet y máximos de coeficientes.

\section{CWT}
Se utiliza un banco de filtros CWT configurable y se extraen medidas de energía,
entropía, concentración y localización tiempo-frecuencia.

\section{S-Transform}
Se implementa una S-Transform discreta y limitada en banda, documentando
explícitamente el vector de frecuencias evaluado.

\section{Clasificación y métricas}
El clasificador es deliberadamente ligero y se usa como instrumento de comparación.
Se reportan accuracy, precision, recall, specificity, macro-F1, weighted-F1 y matriz
de confusión.

\section{Complejidad}
El benchmark reporta tiempo, RTF, throughput y memoria de la representación.

\section{Robustez y estadística}
Se estudian 30, 20, 10, 5, 0 y -5 dB. Los intervalos de confianza se estiman
mediante bootstrap por familia para respetar la dependencia entre variantes de SNR.

\section{Resultados}
\textit{Pendiente de la ejecución final reproducible en MATLAB.}

\end{document}
